\section{Evidence-Centric Robustness: Noisy Retrieval and Distraction}

The first branch of robust RAG studies what enters the context. Early work on irrelevant context established that retrieval errors do not merely lower recall; they can actively redirect generation \cite{irrelevantContext2023}. Noise-aware training and robust RALM work then separated random negatives, hard negatives, counterfactual passages, and noisy retrieved evidence as different stressors rather than one generic perturbation \cite{raat2024,rrag037,rrag091}. Benchmark papers such as RGB and URAG made this decomposition operational by testing negative rejection, information integration, counterfactual evidence, and uncertainty under retrieval noise \cite{rgb2023,rrag001}. These works show that ``noise robustness'' is not a single axis but a family of distinct stressors.

Distraction work focuses on negative marginal utility: a passage is harmful not simply because it is irrelevant, but because it changes the answer, confidence, attribution, or reasoning path. Long-context and context-utility studies show that more context can increase the probability of including plausible but misleading passages \cite{rrag003,rrag039,rrag032}. Utility-oriented retrieval and downstream-aware metrics therefore ask whether a passage helps the generator rather than whether it only matches the query \cite{rrag058,rrag067,rrag073,rrag081}. Context compression and filtering methods try to remove irrelevant spans or documents while preserving evidence needed for multi-hop reasoning \cite{rrag044,chainOfNote2023,sufficientContext2024,rrag043}.

\section{Model-Centric Robustness: Conflict, Faithfulness, and Attribution}

Once evidence enters the prompt, robustness depends on how the generator uses it. Knowledge-conflict papers study contradictions between parametric memory and retrieved evidence, or among retrieved sources \cite{tugOfWar2024,faithfulRAG2025,astuteRAG2025,rrag046}. These papers show that a model can over-trust retrieved evidence, resist correct retrieval because of parametric priors, or silently choose one source without exposing disagreement. Conflict-aware methods therefore localize discrepancies, compare internal and external knowledge, or introduce source-aware reasoning before generation \cite{rrag008,rrag013,rrag017,rrag024,rrag029}.

Faithfulness and attribution papers study whether generated claims are supported by evidence. RAGTruth and related grounding work provide fine-grained hallucination or support labels, while source-preference and entity-conflict studies examine whether models cite and trust the right evidence \cite{ragTruth2024,rrag110,rrag112,rrag115,rrag132}. These papers connect robust RAG to citation verification: document-level citation is often too coarse, because a document can contain multiple claims, qualifiers, and contradictions. Claim-level or sentence-level attribution is more informative but increases annotation and verification cost \cite{rrag076,rrag077,rrag082,rrag086,rrag089}.

\section{Adversarial Evidence Corruption and Security}

Adversarial RAG work asks what happens when evidence unreliability is created by an adaptive attacker. Corpus and retrieval poisoning papers show that small numbers of malicious passages can manipulate dense retrieval or steer final answers \cite{poisoningRetrievalCorpora2023,poisonedRAG2024,badRAG2024,rrag049,rrag057}. Certified and aggregation-based defenses study robustness under bounded corruption assumptions \cite{certifiedRAG2024,rrag016,rrag020}, while reranking and detection defenses use gradient, attention, clustering, or anomaly signals to demote suspicious evidence \cite{rrag004,rrag009,rrag011,rrag015,rrag017}.

Prompt-injection and backdoor work targets a different interface. Instead of inserting false facts, the attacker may insert instructions that the model treats as commands, or triggers that activate malicious behavior through retrieved contexts \cite{aip2025,rrag005,trojanRAG2024,phantom2024,rrag010}. These attacks require retrieval realism: a malicious document must be retrieved under plausible queries, not only under attacker-crafted prompts. This is why we separate factual poisoning, indirect prompt injection, backdoor triggers, and prompt-template attacks as different threat games even when they all appear as retrieved text.

\section{System Control: Calibration, Abstention, and Adaptive Retrieval}

System-control work converts evidence judgments into actions. Answerability and sufficient-context papers ask whether the retrieved context contains enough evidence to answer \cite{sufficientContext2024,rrag111,rrag141,rrag193}. Uncertainty and calibration work asks whether confidence reflects answer correctness, evidence sufficiency, retrieval usefulness, or selective risk \cite{calibRAG2025,uncertaintyRAG2024,rrag001,rrag007,rrag012}. Adaptive retrieval work asks when the system should retrieve, retrieve again, or answer from internal knowledge \cite{selfRAG2023,whenToRetrieve2024,crag2024,ctrla2025,rrag034}.

The control literature highlights a recurring failure: confidence is not the same as evidence sufficiency. A model can be confident in an unsupported answer, or uncertain even when evidence is present but hard to integrate. Trust-or-abstain and self-aware RAG methods therefore model whether to answer, abstain, retrieve more, or ask for clarification \cite{rrag002,rrag006,rrag041,rrag064,rrag066}. Multi-turn and underspecified-query benchmarks add another layer: the correct action may be clarification before retrieval, not direct answering \cite{rrag141,rrag146,rrag153}.

\section{Evaluation and Benchmarks}

Evaluation papers provide the diagnostic infrastructure for robust RAG. RGB decomposes RAG reliability into multiple abilities \cite{rgb2023}; RAGTruth provides fine-grained hallucination annotations \cite{ragTruth2024}; query-robustness and context-utility work stress retrieval interfaces \cite{queryRobustness2025,rrag058}; security benchmarks and arenas evaluate poisoning and defense assumptions \cite{rrag019,rrag023,rrag035}; and uncertainty benchmarks evaluate whether systems know when their answers are reliable \cite{rrag001,rrag093}. Survey papers on RAG, trustworthy RAG, secure RAG, and evaluation provide the broader landscape \cite{ragSurvey2024,ragComprehensive2025,ragEvaluationSurvey2024,trustRAGCompass2024,trustworthyRAGSurvey2025}.

What is still missing is joint evaluation. Existing benchmarks often isolate one failure mode: noise, conflict, poisoning, or unanswerability. Real systems face combinations. A benchmark that includes a stale source, a hard negative, insufficient evidence, and a malicious document will test a different capability from any single-axis benchmark. This motivates the action-trace blueprint: robust-RAG evaluation should score evidence state, answer correctness, action choice, cost, and threat assumptions together.
